% Options for packages loaded elsewhere
\PassOptionsToPackage{unicode}{hyperref}
\PassOptionsToPackage{hyphens}{url}
%
\documentclass[
  ignorenonframetext,
]{beamer}
\usepackage{pgfpages}
\setbeamertemplate{caption}[numbered]
\setbeamertemplate{caption label separator}{: }
\setbeamercolor{caption name}{fg=normal text.fg}
\beamertemplatenavigationsymbolsempty
% Prevent slide breaks in the middle of a paragraph
\widowpenalties 1 10000
\raggedbottom
\setbeamertemplate{part page}{
  \centering
  \begin{beamercolorbox}[sep=16pt,center]{part title}
    \usebeamerfont{part title}\insertpart\par
  \end{beamercolorbox}
}
\setbeamertemplate{section page}{
  \centering
  \begin{beamercolorbox}[sep=12pt,center]{part title}
    \usebeamerfont{section title}\insertsection\par
  \end{beamercolorbox}
}
\setbeamertemplate{subsection page}{
  \centering
  \begin{beamercolorbox}[sep=8pt,center]{part title}
    \usebeamerfont{subsection title}\insertsubsection\par
  \end{beamercolorbox}
}
\AtBeginPart{
  \frame{\partpage}
}
\AtBeginSection{
  \ifbibliography
  \else
    \frame{\sectionpage}
  \fi
}
\AtBeginSubsection{
  \frame{\subsectionpage}
}
\usepackage{amsmath,amssymb}
\usepackage{lmodern}
\usepackage{iftex}
\ifPDFTeX
  \usepackage[T1]{fontenc}
  \usepackage[utf8]{inputenc}
  \usepackage{textcomp} % provide euro and other symbols
\else % if luatex or xetex
  \usepackage{unicode-math}
  \defaultfontfeatures{Scale=MatchLowercase}
  \defaultfontfeatures[\rmfamily]{Ligatures=TeX,Scale=1}
\fi
\usetheme[]{Montpellier}
% Use upquote if available, for straight quotes in verbatim environments
\IfFileExists{upquote.sty}{\usepackage{upquote}}{}
\IfFileExists{microtype.sty}{% use microtype if available
  \usepackage[]{microtype}
  \UseMicrotypeSet[protrusion]{basicmath} % disable protrusion for tt fonts
}{}
\makeatletter
\@ifundefined{KOMAClassName}{% if non-KOMA class
  \IfFileExists{parskip.sty}{%
    \usepackage{parskip}
  }{% else
    \setlength{\parindent}{0pt}
    \setlength{\parskip}{6pt plus 2pt minus 1pt}}
}{% if KOMA class
  \KOMAoptions{parskip=half}}
\makeatother
\usepackage{xcolor}
\IfFileExists{xurl.sty}{\usepackage{xurl}}{} % add URL line breaks if available
\IfFileExists{bookmark.sty}{\usepackage{bookmark}}{\usepackage{hyperref}}
\hypersetup{
  hidelinks,
  pdfcreator={LaTeX via pandoc}}
\urlstyle{same} % disable monospaced font for URLs
\newif\ifbibliography
\usepackage{color}
\usepackage{fancyvrb}
\newcommand{\VerbBar}{|}
\newcommand{\VERB}{\Verb[commandchars=\\\{\}]}
\DefineVerbatimEnvironment{Highlighting}{Verbatim}{commandchars=\\\{\}}
% Add ',fontsize=\small' for more characters per line
\usepackage{framed}
\definecolor{shadecolor}{RGB}{248,248,248}
\newenvironment{Shaded}{\begin{snugshade}}{\end{snugshade}}
\newcommand{\AlertTok}[1]{\textcolor[rgb]{0.94,0.16,0.16}{#1}}
\newcommand{\AnnotationTok}[1]{\textcolor[rgb]{0.56,0.35,0.01}{\textbf{\textit{#1}}}}
\newcommand{\AttributeTok}[1]{\textcolor[rgb]{0.77,0.63,0.00}{#1}}
\newcommand{\BaseNTok}[1]{\textcolor[rgb]{0.00,0.00,0.81}{#1}}
\newcommand{\BuiltInTok}[1]{#1}
\newcommand{\CharTok}[1]{\textcolor[rgb]{0.31,0.60,0.02}{#1}}
\newcommand{\CommentTok}[1]{\textcolor[rgb]{0.56,0.35,0.01}{\textit{#1}}}
\newcommand{\CommentVarTok}[1]{\textcolor[rgb]{0.56,0.35,0.01}{\textbf{\textit{#1}}}}
\newcommand{\ConstantTok}[1]{\textcolor[rgb]{0.00,0.00,0.00}{#1}}
\newcommand{\ControlFlowTok}[1]{\textcolor[rgb]{0.13,0.29,0.53}{\textbf{#1}}}
\newcommand{\DataTypeTok}[1]{\textcolor[rgb]{0.13,0.29,0.53}{#1}}
\newcommand{\DecValTok}[1]{\textcolor[rgb]{0.00,0.00,0.81}{#1}}
\newcommand{\DocumentationTok}[1]{\textcolor[rgb]{0.56,0.35,0.01}{\textbf{\textit{#1}}}}
\newcommand{\ErrorTok}[1]{\textcolor[rgb]{0.64,0.00,0.00}{\textbf{#1}}}
\newcommand{\ExtensionTok}[1]{#1}
\newcommand{\FloatTok}[1]{\textcolor[rgb]{0.00,0.00,0.81}{#1}}
\newcommand{\FunctionTok}[1]{\textcolor[rgb]{0.00,0.00,0.00}{#1}}
\newcommand{\ImportTok}[1]{#1}
\newcommand{\InformationTok}[1]{\textcolor[rgb]{0.56,0.35,0.01}{\textbf{\textit{#1}}}}
\newcommand{\KeywordTok}[1]{\textcolor[rgb]{0.13,0.29,0.53}{\textbf{#1}}}
\newcommand{\NormalTok}[1]{#1}
\newcommand{\OperatorTok}[1]{\textcolor[rgb]{0.81,0.36,0.00}{\textbf{#1}}}
\newcommand{\OtherTok}[1]{\textcolor[rgb]{0.56,0.35,0.01}{#1}}
\newcommand{\PreprocessorTok}[1]{\textcolor[rgb]{0.56,0.35,0.01}{\textit{#1}}}
\newcommand{\RegionMarkerTok}[1]{#1}
\newcommand{\SpecialCharTok}[1]{\textcolor[rgb]{0.00,0.00,0.00}{#1}}
\newcommand{\SpecialStringTok}[1]{\textcolor[rgb]{0.31,0.60,0.02}{#1}}
\newcommand{\StringTok}[1]{\textcolor[rgb]{0.31,0.60,0.02}{#1}}
\newcommand{\VariableTok}[1]{\textcolor[rgb]{0.00,0.00,0.00}{#1}}
\newcommand{\VerbatimStringTok}[1]{\textcolor[rgb]{0.31,0.60,0.02}{#1}}
\newcommand{\WarningTok}[1]{\textcolor[rgb]{0.56,0.35,0.01}{\textbf{\textit{#1}}}}
\setlength{\emergencystretch}{3em} % prevent overfull lines
\providecommand{\tightlist}{%
  \setlength{\itemsep}{0pt}\setlength{\parskip}{0pt}}
\setcounter{secnumdepth}{-\maxdimen} % remove section numbering
%\documentclass[unicode,10pt]{beamer}% 'unicode'が必要
\usepackage{luatexja}% 日本語を使う場合は必須
%\usepackage[japanese]{babel}	
%\usefonttheme[onlymath]{serif}
%\usepackage[T1]{fontenc} %これを使うと、ギリシャ文字が出ない
%\usepackage{textcomp}
%\usepackage[scale=1.0]{tgheros} % Sans serif font
%\usepackage[scaled]{beramono} % Monospace font
%\usepackage{luatexja-otf}
%\renewcommand{\kanjifamilydefault}{\gtdefault}
%\usepackage[haranoaji]{luatexja-preset}

% スライド番号の表示 %%%%%%%%%%%%%%%%%%%
\makeatletter
\setbeamertemplate{footline}{%
  \leavevmode%
  \hbox{%
  \begin{beamercolorbox}[wd=.5\paperwidth,ht=2.25ex,dp=1ex,right]{date in head/foot}%
    \insertframenumber{} \hspace*{2ex} % new version without total frames
  \end{beamercolorbox}}%
  \vskip0pt%
}
\makeatother
% スライド番号の表示 %%%%%%%%%%%%%%%%%%%
\ifLuaTeX
  \usepackage{selnolig}  % disable illegal ligatures
\fi

\title{第6章: 確率}
\subtitle{今井耕介 著\\
『社会科学のためのデータ分析入門 (QSS)』}
\author{}
\date{\vspace{-2.5em}2026-03-09}

\begin{document}
\frame{\titlepage}

\hypertarget{ux78baux7387ux306eux57faux672c}{%
\section{6.1 確率の基本}\label{ux78baux7387ux306eux57faux672c}}

\begin{frame}[fragile]{確率の定義と誕生日問題}
\protect\hypertarget{ux78baux7387ux306eux5b9aux7fa9ux3068ux8a95ux751fux65e5ux554fux984c}{}
\begin{itemize}
\tightlist
\item
  \textbf{頻度主義}: 無限回の試行における相対頻度。
\item
  \textbf{ベイズ主義}: 主観的な確信の度合い。
\end{itemize}

\scriptsize

\begin{Shaded}
\begin{Highlighting}[]
\NormalTok{birthday }\OtherTok{\textless{}{-}} \ControlFlowTok{function}\NormalTok{(k) \{}
\NormalTok{    logdenom }\OtherTok{\textless{}{-}}\NormalTok{ k }\SpecialCharTok{*} \FunctionTok{log}\NormalTok{(}\DecValTok{365}\NormalTok{) }\SpecialCharTok{+} \FunctionTok{lfactorial}\NormalTok{(}\DecValTok{365} \SpecialCharTok{{-}}\NormalTok{ k)}
\NormalTok{    lognumer }\OtherTok{\textless{}{-}} \FunctionTok{lfactorial}\NormalTok{(}\DecValTok{365}\NormalTok{)}
    \DecValTok{1} \SpecialCharTok{{-}} \FunctionTok{exp}\NormalTok{(lognumer }\SpecialCharTok{{-}}\NormalTok{ logdenom)}
\NormalTok{\}}
\NormalTok{k }\OtherTok{\textless{}{-}} \DecValTok{1}\SpecialCharTok{:}\DecValTok{50}\NormalTok{; bday }\OtherTok{\textless{}{-}} \FunctionTok{birthday}\NormalTok{(k)}
\FunctionTok{plot}\NormalTok{(k, bday, }\AttributeTok{xlab =} \StringTok{"Number of people"}\NormalTok{, }\AttributeTok{ylab =} \StringTok{"Probability"}\NormalTok{,}
     \AttributeTok{main =} \StringTok{"Birthday Problem"}\NormalTok{)}
\FunctionTok{abline}\NormalTok{(}\AttributeTok{h =} \FloatTok{0.5}\NormalTok{, }\AttributeTok{lty =} \StringTok{"dashed"}\NormalTok{)}
\end{Highlighting}
\end{Shaded}

\begin{center}\includegraphics[height=0.4\textheight]{QSS_Ch06_files/figure-beamer/unnamed-chunk-1-1} \end{center}
\normalsize
\end{frame}

\hypertarget{ux6761ux4ef6ux4ed8ux304dux78baux7387}{%
\section{6.2 条件付き確率}\label{ux6761ux4ef6ux4ed8ux304dux78baux7387}}

\begin{frame}[fragile]{結合確率・周辺確率・条件付き確率}
\protect\hypertarget{ux7d50ux5408ux78baux7387ux5468ux8fbaux78baux7387ux6761ux4ef6ux4ed8ux304dux78baux7387}{}
\begin{itemize}
\tightlist
\item
  フロリダ州の有権者データ（\texttt{FLVoters.csv}）を使用。
\end{itemize}

\scriptsize

\begin{Shaded}
\begin{Highlighting}[]
\NormalTok{FLVoters }\OtherTok{\textless{}{-}} \FunctionTok{na.omit}\NormalTok{(}\FunctionTok{read.csv}\NormalTok{(}\StringTok{"FLVoters.csv"}\NormalTok{))}
\CommentTok{\# 周辺確率（人種）}
\NormalTok{margin.race }\OtherTok{\textless{}{-}} \FunctionTok{prop.table}\NormalTok{(}\FunctionTok{table}\NormalTok{(FLVoters}\SpecialCharTok{$}\NormalTok{race))}
\CommentTok{\# 条件付き確率（女性における人種の割合）}
\FunctionTok{prop.table}\NormalTok{(}\FunctionTok{table}\NormalTok{(FLVoters}\SpecialCharTok{$}\NormalTok{race[FLVoters}\SpecialCharTok{$}\NormalTok{gender }\SpecialCharTok{==} \StringTok{"f"}\NormalTok{]))}
\end{Highlighting}
\end{Shaded}

\begin{verbatim}
## 
##       asian       black    hispanic      native       other       white 
## 0.016997747 0.138849068 0.136391563 0.003481466 0.032357157 0.671922998
\end{verbatim}

\begin{Shaded}
\begin{Highlighting}[]
\CommentTok{\# 結合確率（人種と性別）}
\NormalTok{joint.p }\OtherTok{\textless{}{-}} \FunctionTok{prop.table}\NormalTok{(}\FunctionTok{table}\NormalTok{(}\AttributeTok{race =}\NormalTok{ FLVoters}\SpecialCharTok{$}\NormalTok{race, }\AttributeTok{gender =}\NormalTok{ FLVoters}\SpecialCharTok{$}\NormalTok{gender))}
\NormalTok{joint.p}
\end{Highlighting}
\end{Shaded}

\begin{verbatim}
##           gender
## race                 f           m
##   asian    0.009107868 0.010095468
##   black    0.074399210 0.056622408
##   hispanic 0.073082410 0.057719741
##   native   0.001865467 0.001316800
##   other    0.017337869 0.016679469
##   white    0.360035115 0.321738176
\end{verbatim}

\normalsize
\end{frame}

\begin{frame}[fragile]{独立性 (Independence)}
\protect\hypertarget{ux72ecux7acbux6027-independence}{}
\begin{itemize}
\tightlist
\item
  \(P(A, B) = P(A)P(B)\) が成り立つとき、\(A\) と \(B\) は独立である。
\end{itemize}

\scriptsize

\begin{Shaded}
\begin{Highlighting}[]
\NormalTok{margin.gender }\OtherTok{\textless{}{-}} \FunctionTok{prop.table}\NormalTok{(}\FunctionTok{table}\NormalTok{(FLVoters}\SpecialCharTok{$}\NormalTok{gender))}
\CommentTok{\# 周辺確率の積 vs 結合確率}
\FunctionTok{plot}\NormalTok{(}\FunctionTok{c}\NormalTok{(margin.race }\SpecialCharTok{*}\NormalTok{ margin.gender[}\StringTok{"f"}\NormalTok{]), }\FunctionTok{c}\NormalTok{(joint.p[, }\StringTok{"f"}\NormalTok{]),}
     \AttributeTok{xlab =} \StringTok{"P(race) * P(female)"}\NormalTok{, }\AttributeTok{ylab =} \StringTok{"P(race and female)"}\NormalTok{,}
     \AttributeTok{main =} \StringTok{"Check for Independence"}\NormalTok{)}
\FunctionTok{abline}\NormalTok{(}\DecValTok{0}\NormalTok{, }\DecValTok{1}\NormalTok{, }\AttributeTok{lty =} \StringTok{"dashed"}\NormalTok{)}
\end{Highlighting}
\end{Shaded}

\begin{center}\includegraphics[height=0.4\textheight]{QSS_Ch06_files/figure-beamer/unnamed-chunk-3-1} \end{center}
\normalsize
\end{frame}

\hypertarget{ux78baux7387ux5206ux5e03}{%
\section{6.3 確率分布}\label{ux78baux7387ux5206ux5e03}}

\begin{frame}[fragile]{二項分布 (Binomial Distribution)}
\protect\hypertarget{ux4e8cux9805ux5206ux5e03-binomial-distribution}{}
\begin{itemize}
\tightlist
\item
  \(n\) 回の試行で成功する回数の分布。
\end{itemize}

\scriptsize

\begin{Shaded}
\begin{Highlighting}[]
\CommentTok{\# n=10, p=0.5 のとき、成功回数が5回となる確率}
\FunctionTok{dbinom}\NormalTok{(}\DecValTok{5}\NormalTok{, }\AttributeTok{size =} \DecValTok{10}\NormalTok{, }\AttributeTok{prob =} \FloatTok{0.5}\NormalTok{)}
\end{Highlighting}
\end{Shaded}

\begin{verbatim}
## [1] 0.2460938
\end{verbatim}

\begin{Shaded}
\begin{Highlighting}[]
\CommentTok{\# 累積分布関数: 成功回数が3回以下となる確率}
\FunctionTok{pbinom}\NormalTok{(}\DecValTok{3}\NormalTok{, }\AttributeTok{size =} \DecValTok{10}\NormalTok{, }\AttributeTok{prob =} \FloatTok{0.5}\NormalTok{)}
\end{Highlighting}
\end{Shaded}

\begin{verbatim}
## [1] 0.171875
\end{verbatim}

\normalsize
\end{frame}

\begin{frame}[fragile]{正規分布 (Normal Distribution)}
\protect\hypertarget{ux6b63ux898fux5206ux5e03-normal-distribution}{}
\begin{itemize}
\tightlist
\item
  連続変数の最も重要な分布。
\end{itemize}

\scriptsize

\begin{Shaded}
\begin{Highlighting}[]
\CommentTok{\# 標準正規分布において、平均から ±1 標準偏差の範囲に入る確率}
\FunctionTok{pnorm}\NormalTok{(}\DecValTok{1}\NormalTok{) }\SpecialCharTok{{-}} \FunctionTok{pnorm}\NormalTok{(}\SpecialCharTok{{-}}\DecValTok{1}\NormalTok{)}
\end{Highlighting}
\end{Shaded}

\begin{verbatim}
## [1] 0.6826895
\end{verbatim}

\begin{Shaded}
\begin{Highlighting}[]
\CommentTok{\# ±2 標準偏差の範囲（約 95\%）}
\FunctionTok{pnorm}\NormalTok{(}\DecValTok{2}\NormalTok{) }\SpecialCharTok{{-}} \FunctionTok{pnorm}\NormalTok{(}\SpecialCharTok{{-}}\DecValTok{2}\NormalTok{)}
\end{Highlighting}
\end{Shaded}

\begin{verbatim}
## [1] 0.9544997
\end{verbatim}

\normalsize
\end{frame}

\hypertarget{ux5927ux6a19ux672cux5b9aux7406}{%
\section{6.4 大標本定理}\label{ux5927ux6a19ux672cux5b9aux7406}}

\begin{frame}[fragile]{大数の法則 (Law of Large Numbers)}
\protect\hypertarget{ux5927ux6570ux306eux6cd5ux5247-law-of-large-numbers}{}
\begin{itemize}
\tightlist
\item
  サンプルサイズが大きくなるほど、標本平均は期待値に収束する。
\end{itemize}

\scriptsize

\begin{Shaded}
\begin{Highlighting}[]
\NormalTok{sims }\OtherTok{\textless{}{-}} \DecValTok{1000}
\NormalTok{x.unif }\OtherTok{\textless{}{-}} \FunctionTok{runif}\NormalTok{(sims) }\CommentTok{\# 一様分布 [0, 1]}
\NormalTok{mean.unif }\OtherTok{\textless{}{-}} \FunctionTok{cumsum}\NormalTok{(x.unif) }\SpecialCharTok{/} \DecValTok{1}\SpecialCharTok{:}\NormalTok{sims}
\FunctionTok{plot}\NormalTok{(}\DecValTok{1}\SpecialCharTok{:}\NormalTok{sims, mean.unif, }\AttributeTok{type =} \StringTok{"l"}\NormalTok{, }\AttributeTok{xlab =} \StringTok{"Sample size"}\NormalTok{, }
     \AttributeTok{ylab =} \StringTok{"Sample mean"}\NormalTok{, }\AttributeTok{main =} \StringTok{"Law of Large Numbers"}\NormalTok{)}
\FunctionTok{abline}\NormalTok{(}\AttributeTok{h =} \FloatTok{0.5}\NormalTok{, }\AttributeTok{lty =} \StringTok{"dashed"}\NormalTok{)}
\end{Highlighting}
\end{Shaded}

\begin{center}\includegraphics[height=0.4\textheight]{QSS_Ch06_files/figure-beamer/unnamed-chunk-6-1} \end{center}
\normalsize
\end{frame}

\begin{frame}[fragile]{中心極限定理 (Central Limit Theorem)}
\protect\hypertarget{ux4e2dux5fc3ux6975ux9650ux5b9aux7406-central-limit-theorem}{}
\begin{itemize}
\tightlist
\item
  サンプルサイズが大きければ、標本平均の分布は正規分布に従う。
\end{itemize}

\scriptsize

\begin{Shaded}
\begin{Highlighting}[]
\NormalTok{z.unif }\OtherTok{\textless{}{-}} \FunctionTok{rep}\NormalTok{(}\ConstantTok{NA}\NormalTok{, }\DecValTok{1000}\NormalTok{)}
\ControlFlowTok{for}\NormalTok{ (i }\ControlFlowTok{in} \DecValTok{1}\SpecialCharTok{:}\DecValTok{1000}\NormalTok{) \{}
\NormalTok{    x }\OtherTok{\textless{}{-}} \FunctionTok{runif}\NormalTok{(}\DecValTok{100}\NormalTok{, }\AttributeTok{min =} \DecValTok{0}\NormalTok{, }\AttributeTok{max =} \DecValTok{1}\NormalTok{) }\CommentTok{\# n=100}
\NormalTok{    z.unif[i] }\OtherTok{\textless{}{-}}\NormalTok{ (}\FunctionTok{mean}\NormalTok{(x) }\SpecialCharTok{{-}} \FloatTok{0.5}\NormalTok{) }\SpecialCharTok{/} \FunctionTok{sqrt}\NormalTok{(}\DecValTok{1} \SpecialCharTok{/}\NormalTok{ (}\DecValTok{12} \SpecialCharTok{*} \DecValTok{100}\NormalTok{))}
\NormalTok{\}}
\FunctionTok{hist}\NormalTok{(z.unif, }\AttributeTok{freq =} \ConstantTok{FALSE}\NormalTok{, }\AttributeTok{nclass =} \DecValTok{30}\NormalTok{, }\AttributeTok{xlab =} \StringTok{"z{-}score"}\NormalTok{, }
     \AttributeTok{main =} \StringTok{"Central Limit Theorem (Uniform)"}\NormalTok{)}
\FunctionTok{curve}\NormalTok{(}\FunctionTok{dnorm}\NormalTok{(x), }\AttributeTok{add =} \ConstantTok{TRUE}\NormalTok{, }\AttributeTok{col =} \StringTok{"red"}\NormalTok{)}
\end{Highlighting}
\end{Shaded}

\begin{center}\includegraphics[height=0.4\textheight]{QSS_Ch06_files/figure-beamer/unnamed-chunk-7-1} \end{center}
\normalsize
\end{frame}

\hypertarget{ux307eux3068ux3081}{%
\section{6.5 まとめ}\label{ux307eux3068ux3081}}

\begin{frame}[fragile]{第6章のまとめ}
\protect\hypertarget{ux7b2c6ux7ae0ux306eux307eux3068ux3081}{}
\begin{itemize}
\tightlist
\item
  \textbf{確率}: 不確実な出来事を記述するための数学的な枠組み。
\item
  \textbf{条件付き確率と独立性}: 変数間の関係を理解するための鍵。
\item
  \textbf{確率分布}:
  データの発生メカニズムをモデル化する（二項分布、正規分布など）。
\item
  \textbf{大標本定理}:
  標本から母集団を推測するための理論的基礎（大数の法則、中心極限定理）。
\item
  Rの操作: \texttt{sample()}, \texttt{dbinom()}, \texttt{pnorm()},
  \texttt{runif()}, \texttt{rbinom()} などを習得した。
\end{itemize}
\end{frame}

\end{document}
